\documentclass[10pt,a4paper]{amsart}
\usepackage{amsmath,amssymb,amsthm,graphicx}
% Retained minimal preamble for the Stacks Project chapter more-algebra.tex (More on Algebra).
% The chapter file pinned-original-source.tex carries only \input{preamble} on its first line, so no
% preamble text lives inside the chapter itself.  The Stacks Project uses one single shared preamble for
% every chapter; the snapshot stages it identically as sources/stacks-chow/preamble.tex and
% sources/stacks-etale-cohomology/preamble.tex.  The block below is copied mechanically, line for line,
% from sources/stacks-chow/preamble.tex lines 14-18 (the Xy-pic setup) and 154-217 (the native theorem
% declarations together with the mathematical macros).  Nothing was added, removed or redefined; the
% wrapper class and the amsmath/amssymb/amsthm/graphicx base packages are supplied by the assembler.
% Chapter: more-algebra.tex -- More on Algebra.
% For commutative diagrams we use Xy-pic
\usepackage[all]{xy}

% We use 2cell for 2-commutative diagrams.
\xyoption{2cell}

% Theorem environments.
%
\theoremstyle{plain}
\newtheorem{theorem}[subsection]{Theorem}
\newtheorem{proposition}[subsection]{Proposition}
\newtheorem{lemma}[subsection]{Lemma}

\theoremstyle{definition}
\newtheorem{definition}[subsection]{Definition}
\newtheorem{example}[subsection]{Example}
\newtheorem{exercise}[subsection]{Exercise}
\newtheorem{situation}[subsection]{Situation}

\theoremstyle{remark}
\newtheorem{remark}[subsection]{Remark}
\newtheorem{remarks}[subsection]{Remarks}

\numberwithin{equation}{subsection}

% Macros
%
\def\lim{\mathop{\mathrm{lim}}\nolimits}
\def\colim{\mathop{\mathrm{colim}}\nolimits}
\def\Spec{\mathop{\mathrm{Spec}}}
\def\Hom{\mathop{\mathrm{Hom}}\nolimits}
\def\Ext{\mathop{\mathrm{Ext}}\nolimits}
\def\SheafHom{\mathop{\mathcal{H}\!\mathit{om}}\nolimits}
\def\SheafExt{\mathop{\mathcal{E}\!\mathit{xt}}\nolimits}
\def\Sch{\mathit{Sch}}
\def\Mor{\mathop{\mathrm{Mor}}\nolimits}
\def\Ob{\mathop{\mathrm{Ob}}\nolimits}
\def\Sh{\mathop{\mathit{Sh}}\nolimits}
\def\NL{\mathop{N\!L}\nolimits}
\def\CH{\mathop{\mathrm{CH}}\nolimits}
\def\proetale{{pro\text{-}\acute{e}tale}}
\def\etale{{\acute{e}tale}}
\def\QCoh{\mathit{QCoh}}
\def\Ker{\mathop{\mathrm{Ker}}}
\def\Im{\mathop{\mathrm{Im}}}
\def\Coker{\mathop{\mathrm{Coker}}}
\def\Coim{\mathop{\mathrm{Coim}}}

% Boxtimes
%
\DeclareMathSymbol{\boxtimes}{\mathbin}{AMSa}{"02}

%
% Macros for moduli stacks/spaces
%
\def\QCohstack{\mathcal{QC}\!\mathit{oh}}
\def\Cohstack{\mathcal{C}\!\mathit{oh}}
\def\Spacesstack{\mathcal{S}\!\mathit{paces}}
\def\Quotfunctor{\mathrm{Quot}}
\def\Hilbfunctor{\mathrm{Hilb}}
\def\Curvesstack{\mathcal{C}\!\mathit{urves}}
\def\Polarizedstack{\mathcal{P}\!\mathit{olarized}}
\def\Complexesstack{\mathcal{C}\!\mathit{omplexes}}
% \Pic is the operator that assigns to X its picard group, usage \Pic(X)
% \Picardstack_{X/B} denotes the Picard stack of X over B
% \Picardfunctor_{X/B} denotes the Picard functor of X over B
\def\Pic{\mathop{\mathrm{Pic}}\nolimits}
\def\Picardstack{\mathcal{P}\!\mathit{ic}}
\def\Picardfunctor{\mathrm{Pic}}
\def\Deformationcategory{\mathcal{D}\!\mathit{ef}}
\title{Multiplication by a power of f identifies the cohomology of M up to f-power torsion with the cohomology of eta\_f M}
\author{The Stacks Project Authors}
\date{stacks-project:more-algebra.tex:local-snapshot-20261012; GNU FDL 1.2 or later}
\begin{document}
\maketitle
\setcounter{subsection}{1}
\noindent\textit{Editorial note.} All delivered mathematical text is the author's own, copied from the pinned Stacks Project chapter file more-algebra.tex; no statement or proof step was rewritten, reordered or removed. The counted statement is the labeled lemma environment at source lines 27493-27501 (canonical tag 0F7P) and the counted proof is the proof environment at source lines 27503-27515, which immediately follows it and is its own proof.\par\medskip
\noindent\textit{Editorial note.} The notation eta\_f M of the statement and proof denotes the shifted complex with i-th term f\textasciicircum{}i M\textasciicircum{}i, and it is written by the author with the ordinary Greek letter eta and the subscript f; no macro beyond the shared preamble is needed, and the example that precedes the lemma at source lines 27480-27491 is outside the counted spans.\par\medskip
\noindent\textit{Editorial note.} The chapter file pinned-original-source.tex carries only an input of the shared preamble on its first line, so no preamble text lives inside the chapter. The retained file sources/stacks-more-algebra/retained-preamble.tex is a mechanical, line-for-line copy of the project's single shared preamble as staged in the snapshot (Xy-pic setup plus the native theorem declarations and mathematical macros). Nothing was added, removed or redefined, and the excerpt serves only as preamble, never as a span coordinate.\par\medskip
\noindent\textit{Editorial note.} The source URL is the source-level Stacks Project URL, as required when more than one card is produced from one source; the per-claim tag is carried by claim\_key and by this note. The counted statement and proof refer to no labelled item of the chapter and cite no bibliography entry, so this card contributes no external reference stub and needs no bibliography.\par\medskip
\noindent\textit{Allowed standard prerequisites (editorial).}\begin{itemize}
\item The shifted complex of a complex M with respect to a nonzerodivisor f has i-th term f\textasciicircum{}i M\textasciicircum{}i and differentials induced by the differentials of M. used at proof lines 27504-27506, where the cycles of the shifted complex are computed. Turns the cohomology of the shifted complex into a quotient of f\textasciicircum{}i times the cycles of M.
\item The kernel of the differential in degree i of the shifted complex equals f\textasciicircum{}i ker(d\textasciicircum{}i), because multiplication by f\textasciicircum{}i is injective on M\textasciicircum{}i and on M\textasciicircum{}\{i+1\} when f is a nonzerodivisor on all terms. used at proof lines 27505-27506. Shows that multiplication by f\textasciicircum{}i lands in the cycles of the shifted complex and induces a surjection on cohomology.
\item An element of the kernel of the induced map on cohomology satisfies f\textasciicircum{}i z = d(f\textasciicircum{}\{i-1\} y), which after cancelling f\textasciicircum{}i and f\textasciicircum{}\{i-1\} means f z = d(y), that is z is f-torsion in cohomology. used at proof lines 27510-27514. Identifies the kernel of the surjection with the f-torsion subgroup of H\textasciicircum{}i(M), which yields the displayed isomorphism.
\end{itemize}\medskip
\begin{lemma}
\label{lemma-eta-first-property}
Let $A$ be a ring and let $f \in A$ be a nonzerodivisor. Let $M^\bullet$ be
a complex of $f$-torsion free $A$-modules. There is a canonical isomorphism
$$
f^i : H^i(M^\bullet)/H^i(M^\bullet)[f] \longrightarrow H^i(\eta_fM^\bullet)
$$
given by multiplication by $f^i$.
\end{lemma}

\begin{proof}
Observe that $\Ker(d^i : (\eta_fM)^i \to (\eta_fM)^{i + 1})$
is equal to
$\Ker(d^i : f^iM^i \to f^iM^{i + 1}) = f^i\Ker(d^i : M^i \to M^{i + 1})$.
This we get a surjection
$f^i : H^i(M^\bullet) \to H^i(\eta_fM^\bullet)$
by sending the class of $z \in \Ker(d^i : M^i \to M^{i + 1})$
to the class of $f^iz$. If we obtain the zero class in
$H^i(\eta_fM^\bullet)$ then we see that $f^i z = d^{i - 1}(f^{i - 1}y)$
for some $y \in M^{i - 1}$. Since $f$ is a nonzerodivisor on all
the modules involved, this means $f z = d^{i - 1}(y)$ which
exactly means that the class of $z$ is $f$-torsion as desired.
\end{proof}

\end{document}
